% GENERATED FILE. Edit canonical lessons, then run npm run generate:books.
% canonical-source-hash: fnv1a64:92fa48db9b2b716c
% canonical-lessons: PT-C420-boia, PT-C420-ancora, PT-C420-atracar, PT-C420-enjoo, PT-C420-enjoar

\chapter{Chegar a terra}
\label{ch:pt-chegar-a-terra}
\hlchaptermodality{420}

\begin{chapteropening}
\textbf{By the end of this chapter:} \emph{I can talk about reaching land: the buoy, the anchor, mooring, seasickness and feeling sick.}

Complete the last lesson of chapter 420: I can talk about reaching land: the buoy, the anchor, mooring, seasickness and feeling sick.
\end{chapteropening}

\section[a boia]{a boia — the buoy; the float, the rubber ring}

\label{lesson:PT-C420-boia}



\begin{quote}
Before the new one: say the Portuguese for the bow, then the Portuguese for the stern.
\end{quote}

\subsection*{You'll want to know: a boia}
\textbf{a boia} — \textquotedblleft{}the buoy; the float, the rubber ring\textquotedblright{}.

\begin{quote}
As crianças nadam com boias nos braços.
\end{quote}

\begin{quote}
\emph{The children swim with armbands on.}
\end{quote}

\emph{Boias} on the arms are armbands; \emph{uma boia} in the sea marks a safe area.

\begin{cousinweb}
\emph{boia} was borrowed from Old French; EtymDB traces the French word to a Frankish word related to \emph{beacon} 'sign'.
\end{cousinweb}

\begin{grammarlens}[title={What you've built}]
One more word for telling, explaining and describing at B1.
\end{grammarlens}

\subsection*{Your turn}
\begin{itemize}
\raggedright
  \item \emph{Say it:} \emph{a boia}
  \item \emph{Say it:} \emph{a boia}, once more
  \item \emph{Say it:} a sentence of your own with \emph{a boia}
  \item \emph{Recall:} say the Portuguese for the bow, then the Portuguese for the stern, then say \emph{a boia} again
\end{itemize}

\begin{tcolorbox}[breakable,colback=teal!4,colframe=teal!35!black,title={Before you move on}]
What does a boia mean? (\textquotedblleft{}The buoy; the float, the rubber ring\textquotedblright{}.) Say it once more.
\end{tcolorbox}
\section[a âncora]{a âncora — the anchor}

\label{lesson:PT-C420-ancora}



\begin{quote}
Before the new one: say the Portuguese for the stern, then the Portuguese for the buoy.
\end{quote}

\subsection*{You'll want to know: a âncora}
\textbf{a âncora} — \textquotedblleft{}the anchor\textquotedblright{}.

\begin{quote}
O barco lançou a âncora perto da praia.
\end{quote}

\begin{quote}
\emph{The boat dropped anchor near the beach.}
\end{quote}

In Brazil, \emph{o âncora} is also a TV news anchor; Portugal says \emph{o pivô}.

\begin{cousinweb}
\emph{âncora} is inherited from Latin \emph{ancora}, which came from Greek \emph{ánkyra} 'anchor' (EtymDB).
\end{cousinweb}

\begin{grammarlens}[title={What you've built}]
One more word for telling, explaining and describing at B1.
\end{grammarlens}

\subsection*{Your turn}
\begin{itemize}
\raggedright
  \item \emph{Say it:} \emph{a âncora}
  \item \emph{Say it:} \emph{a âncora}, once more
  \item \emph{Say it:} a sentence of your own with \emph{a âncora}
  \item \emph{Recall:} say the Portuguese for the stern, then the Portuguese for the buoy, then say \emph{a âncora} again
\end{itemize}

\begin{tcolorbox}[breakable,colback=teal!4,colframe=teal!35!black,title={Before you move on}]
What does a âncora mean? (\textquotedblleft{}The anchor\textquotedblright{}.) Say it once more.
\end{tcolorbox}
\section[atracar]{atracar — to moor, to dock}

\label{lesson:PT-C420-atracar}



\begin{quote}
Before the new one: say the Portuguese for the buoy, then the Portuguese for the anchor.
\end{quote}

\subsection*{You'll want to know: atracar}
\textbf{atracar} — \textquotedblleft{}to moor, to dock\textquotedblright{}.

\begin{quote}
O ferry atracou no cais às dez.
\end{quote}

\begin{quote}
\emph{The ferry docked at the quay at ten.}
\end{quote}

The opposite is \emph{desatracar}, to cast off.

\begin{cousinweb}
The origin of \emph{atracar} is not settled. EtymDB links it with an Arabic verb meaning 'to advance', but that link is not settled.
\end{cousinweb}

\begin{grammarlens}[title={What you've built}]
One more word for telling, explaining and describing at B1.
\end{grammarlens}

\subsection*{Your turn}
\begin{itemize}
\raggedright
  \item \emph{Say it:} \emph{atracar}
  \item \emph{Say it:} \emph{atracar}, once more
  \item \emph{Say it:} a sentence of your own with \emph{atracar}
  \item \emph{Recall:} say the Portuguese for the buoy, then the Portuguese for the anchor, then say \emph{atracar} again
\end{itemize}

\begin{tcolorbox}[breakable,colback=teal!4,colframe=teal!35!black,title={Before you move on}]
What does atracar mean? (\textquotedblleft{}To moor, to dock\textquotedblright{}.) Say it once more.
\end{tcolorbox}
\section[o enjoo]{o enjoo — seasickness, nausea}

\label{lesson:PT-C420-enjoo}



\begin{quote}
Before the new one: say the Portuguese for the anchor, then the Portuguese for to moor.
\end{quote}

\subsection*{You'll want to know: o enjoo}
\textbf{o enjoo} — \textquotedblleft{}seasickness, nausea\textquotedblright{}.

\begin{quote}
Tomei um comprimido para o enjoo.
\end{quote}

\begin{quote}
\emph{I took a pill for the seasickness.}
\end{quote}

Also car sickness or morning sickness.

\begin{cousinweb}
\emph{enjoo} is built on \emph{enjoar} 'to feel sick', whose own origin is not settled (EtymDB lists \emph{enjoar} without a source).
\end{cousinweb}

\begin{grammarlens}[title={What you've built}]
One more word for telling, explaining and describing at B1.
\end{grammarlens}

\subsection*{Your turn}
\begin{itemize}
\raggedright
  \item \emph{Say it:} \emph{o enjoo}
  \item \emph{Say it:} \emph{o enjoo}, once more
  \item \emph{Say it:} a sentence of your own with \emph{o enjoo}
  \item \emph{Recall:} say the Portuguese for the anchor, then the Portuguese for to moor, then say \emph{o enjoo} again
\end{itemize}

\begin{tcolorbox}[breakable,colback=teal!4,colframe=teal!35!black,title={Before you move on}]
What does o enjoo mean? (\textquotedblleft{}Seasickness, nausea\textquotedblright{}.) Say it once more.
\end{tcolorbox}
\section[enjoar]{enjoar — to feel sick (on a boat or in a car); to get sick of}

\label{lesson:PT-C420-enjoar}



\begin{quote}
Before the new one: say the Portuguese for to moor, then the Portuguese for seasickness.
\end{quote}

\subsection*{You'll want to know: enjoar}
\textbf{enjoar} — \textquotedblleft{}to feel sick (on a boat or in a car); to get sick of\textquotedblright{}.

\begin{quote}
Enjoo sempre que viajo de barco.
\end{quote}

\begin{quote}
\emph{I always feel sick when I travel by boat.}
\end{quote}

Informal: \emph{Já enjoei esta música} — I am sick of this song.

\begin{cousinweb}
\emph{enjoar} is usually traced to Latin \emph{in odiō} 'in hatred', the phrase behind Spanish \emph{enojar}, but the origin is not settled; EtymDB has no source for it.
\end{cousinweb}

\begin{grammarlens}[title={What you've built}]
One more word for telling, explaining and describing at B1.
\end{grammarlens}

\subsection*{Your turn}
\begin{itemize}
\raggedright
  \item \emph{Say it:} \emph{enjoar}
  \item \emph{Say it:} \emph{enjoar}, once more
  \item \emph{Say it:} a sentence of your own with \emph{enjoar}
  \item \emph{Recall:} say the Portuguese for to moor, then the Portuguese for seasickness, then say \emph{enjoar} again
\end{itemize}

\begin{tcolorbox}[breakable,colback=teal!4,colframe=teal!35!black,title={Before you move on}]
What does enjoar mean? (\textquotedblleft{}To feel sick (on a boat or in a car); to get sick of\textquotedblright{}.) Say it once more.
\end{tcolorbox}
